\documentclass[../../main.tex]{subfiles}
\graphicspath{{../../image/}{../../figure/}}
\begin{document}

\section{Cauchy 问题：古典解}\label{section:ch04-经典解}

为避免平凡的复杂情况，我们把研究集中于正则形式的齐次双曲守恒律组
\begin{align}
\partial_t U(x,t)+\operatorname{div}G(U(x,t))=0,
\label{align:ch04-守恒律组}
\end{align}
尽管分析可以按常规方式推广到一般双曲平衡律组\cref{chapter:ch03-双曲平衡律方程组}。空间变量 $x$ 取值于 $\mathbb{R}^{m}$，时间 $t$ 取值于 $[0,T)$，其中 $T>0$，也可能 $T=\infty$。状态向量 $U$ 取值于 $\mathbb{R}^{n}$ 的某个开子集 $\mathcal{O}$，而 $G=(G_{1},\cdots,G_{m})$ 是定义于 $\mathcal{O}$ 上的给定光滑函数，取值于 $\mathbb{M}^{n\times m}$。称\eqref{align:ch04-守恒律组}是双曲的，如果对每一固定的 $U\in\mathcal{O}$ 与 $\nu\in\mathbb{S}^{m-1}$，$n\times n$ 矩阵
\begin{align}
\Lambda(\nu;U)=\sum_{\alpha=1}^{m}\nu_{\alpha}\mathrm{D}G_{\alpha}(U)
\label{align:ch04-方向矩阵}
\end{align}
具有实特征值 $\lambda_{1}(\nu;U),\cdots,\lambda_{n}(\nu;U)$ 与 $n$ 个线性无关的特征向量 $R_{1}(\nu;U),\cdots,R_{n}(\nu;U)$。

为表述 Cauchy 问题，我们附加初始条件
\begin{align}
U(x,0)=U_{0}(x),\qquad x\in\mathbb{R}^{m},
\label{align:ch04-Cauchy初值}
\end{align}
其中 $U_{0}$ 是从 $\mathbb{R}^{m}$ 到 $\mathcal{O}$ 的函数。

\eqref{align:ch04-守恒律组}的\emph{古典解}是指一个局部 Lipschitz 函数 $U$，定义在 $\mathbb{R}^{m}\times[0,T)$ 上、取值于 $\mathcal{O}$，且几乎处处满足\eqref{align:ch04-守恒律组}。若此函数还对所有 $x\in\mathbb{R}^{m}$ 满足\eqref{align:ch04-Cauchy初值}，则称它是以 $U_{0}$ 为初值的 Cauchy 问题的解。

我们将会看到，当方程组配有一个熵 $\eta$ 及相应熵流 $Q$ 时，Cauchy 问题的理论将大为丰富；二者由
\begin{align}
\mathrm{D}Q_{\alpha}(U)=\mathrm{D}\eta(U)\mathrm{D}G_{\alpha}(U),\qquad \alpha=1,\cdots,m
\label{align:ch04-熵流关系}
\end{align}
相联系。此时，\eqref{align:ch04-守恒律组}的任一古典解还满足附加守恒律
\begin{align}
\partial_{t}\eta(U(x,t))+\operatorname{div}Q(U(x,t))=0.
\label{align:ch04-熵守恒}
\end{align}

随着理论的展开，我们将看到凸熵函数对解施加稳定化影响。作为这一效应的第一个迹象，下面的命题表明：对于配有凸熵的方程组，初值对 Cauchy 问题解的影响范围是有限的。

\begin{theorem}[有限影响区域]\label{theorem:ch04-影响区域}
假定\eqref{align:ch04-守恒律组}是双曲组，特征速度 $\lambda_{1}(\nu;U)\leqslant\cdots\leqslant\lambda_{n}(\nu;U)$，并配有熵 $\eta(U)$ 及相应熵流 $Q(U)$。设 $U(x,t)$ 是\eqref{align:ch04-守恒律组}在 $\mathbb{R}^{m}\times[0,T)$ 上以\eqref{align:ch04-Cauchy初值}为初值的古典解，其中 $U_{0}$ 在一个半空间上为常数：存在 $\xi\in\mathbb{S}^{m-1}$，使得当 $x\cdot\xi\geqslant 0$ 时 $U_{0}(x)=\bar{U}=\text{常数}$。再假定 $\mathrm{D}^{2}\eta(\bar{U})$ 是正定的。那么，对任一 $t\in[0,T)$，当 $x\cdot\xi\geqslant\lambda_{n}(\xi;\bar{U})t$ 时 $U(x,t)=\bar{U}$。
\end{theorem}

\begin{proof}
不失一般性，可设 $\eta(\bar{U})=0$，$\mathrm{D}\eta(\bar{U})=0$，$Q_{\alpha}(\bar{U})=0$，$\mathrm{D}Q_{\alpha}(\bar{U})=0$，$\alpha=1,\cdots,m$；否则只需把给定的熵--熵流对替换为新的一对
\begin{align}
\bar{\eta}(U)=\eta(U)-\eta(\bar{U})-\mathrm{D}\eta(\bar{U})[U-\bar{U}],
\label{align:ch04-eta规范化}
\end{align}
\begin{align}
\bar{Q}(U)=Q(U)-Q(\bar{U})-\mathrm{D}\eta(\bar{U})[G(U)-G(\bar{U})].
\label{align:ch04-Q规范化}
\end{align}

对每一 $s\in\mathbb{R}$，$\nu\in\mathbb{S}^{m-1}$，$U\in\mathcal{O}$，定义
\begin{align}
\Phi(s,\nu;U)=s\eta(U)-Q(U)\nu,
\label{align:ch04-Phi}
\end{align}
注意 $\Phi(s,\nu;\bar{U})=0$ 且 $\mathrm{D}\Phi(s,\nu;\bar{U})=0$。进而，利用\eqref{align:ch04-熵流关系}与\eqref{align:ch04-方向矩阵}，
\begin{align}
\mathrm{D}^{2}\Phi(s,\nu;\bar{U})=\mathrm{D}^{2}\eta(\bar{U})[sI-\Lambda(\nu;\bar{U})].
\label{align:ch04-D2Phi}
\end{align}
于是，对 $j,k=1,\cdots,n$，
\begin{align}
R_{j}(\nu;\bar{U})^{\top}\mathrm{D}^{2}\Phi(s,\nu;\bar{U})R_{k}(\nu;\bar{U})=[s-\lambda_{k}(\nu;\bar{U})]R_{j}(\nu;\bar{U})^{\top}\mathrm{D}^{2}\eta(\bar{U})R_{k}(\nu;\bar{U}).
\label{align:ch04-RD2}
\end{align}
由\cref{chapter:ch03-双曲平衡律方程组}，当 $j\neq k$ 时\eqref{align:ch04-RD2}右端为零；当 $j=k$ 时，由于 $\mathrm{D}^{2}\eta(\bar{U})$ 正定，其符号与 $s-\lambda_{k}(\nu;\bar{U})$ 相同。


对 $\varepsilon>0$，令 $\bar{s}=\max_{\nu\in\mathbb{S}^{m-1}}\lambda_{n}(\nu;\bar{U})+\varepsilon$，$\hat{s}=\lambda_{n}(\xi;\bar{U})+\varepsilon$。则存在 $\delta=\delta(\varepsilon)$ 使得
\begin{align}
\begin{cases}
\Phi(\bar{s},\nu;U)>0,&\text{for }0<|U-\bar{U}|<\delta(\varepsilon),\ \nu\in\mathbb{S}^{m-1}\\
\Phi(\hat{s},\xi;U)>0,&\text{for }0<|U-\bar{U}|<\delta(\varepsilon).
\end{cases}
\label{align:ch04-Phi正}
\end{align}

为证明定理的断言，只需证明：对每一固定的 $\varepsilon>0$ 与 $t\in[0,T)$，当 $x\cdot\xi\geqslant\hat{s}t$ 时 $U(x,t)=\bar{U}$。

对任一点 $(y,\tau)$，其中 $\tau\in(0,T)$ 且 $y\cdot\xi\geqslant\hat{s}\tau$，我们关联以锥
\begin{align}
\mathscr{K}_{y,\tau}=\{(x,t):0\leqslant t\leqslant\tau,\ |x-y|\leqslant\bar{s}(\tau-t),\ x\cdot\xi\geqslant y\cdot\xi-\hat{s}(\tau-t)\},
\label{align:ch04-影响锥}
\end{align}
它包含于集合 $\{(x,t):0\leqslant t<T,\ x\cdot\xi\geqslant\hat{s}t\}$。于是 $\mathscr{K}_{y,\tau}$ 的 $t$-截面的边界是超平面中垂直于 $\xi$ 的子集 $\mathscr{P}_{t}$ 与以 $y$ 为心、半径为 $\bar{s}(\tau-t)$ 的球面的子集 $\mathscr{S}_{t}$ 的并 $\mathscr{P}_{t}\cup\mathscr{S}_{t}$。当 $x\in\mathscr{P}_{t}$ 时，$t$-截面在点 $x$ 处的单位外法向为 $-\xi$；当 $x\in\mathscr{S}_{t}$ 时为 $\bar{s}^{-1}(\tau-t)^{-1}(x-y)$。因此，在 $\mathscr{K}_{y,\tau}$ 上积分\eqref{align:ch04-熵守恒}，应用 Green 定理并采用记号\eqref{align:ch04-Phi}，得到
\begin{align}
\int_{0}^{\tau}\int_{\mathscr{P}_{t}}\Phi(\hat{s},\xi;U)\,\mathrm{d}\mathscr{H}^{m-1}(x)\mathrm{d}t+\int_{0}^{\tau}\int_{\mathscr{S}_{t}}\Phi(\bar{s},-\bar{s}^{-1}(\tau-t)^{-1}(x-y);U)\,\mathrm{d}\mathscr{H}^{m-1}(x)\mathrm{d}t=0.
\label{align:ch04-锥积分}
\end{align}

作此准备后，假设定理的断言不成立。由于 $U(x,t)$ 连续，且当 $x\cdot\xi\geqslant 0$ 时 $U(x,0)=\bar{U}$，可找到 $(y,\tau)$，其中 $\tau\in(0,T)$ 且 $y\cdot\xi\geqslant\hat{s}\tau$，使得 $U(y,\tau)\neq\bar{U}$，并且对一切 $(x,t)\in\mathscr{K}_{y,\tau}$ 有 $|U(x,t)-\bar{U}|<\delta(\varepsilon)$。此时，\eqref{align:ch04-锥积分}与\eqref{align:ch04-Phi正}结合即导出矛盾。证明完成。
\end{proof}

有趣的是，在上述证明中，一个粗糙的“能量”估计给出了初值影响范围增长率的精确值。

正如我们将在第 V 章中看到的，只要 $U_{0}$ 充分光滑且 $T$ 充分小，古典解类中的 Cauchy 问题是适定的。然而，整体情形则截然不同。这将在下一节中得到展示。


\end{document}
